\documentclass[11pt,fontset=fandol]{ctexart}
\usepackage[a4paper,margin=23mm]{geometry}
\usepackage{amsmath,amssymb,mathtools,bm,mathrsfs,esint}
\usepackage{graphicx,longtable,booktabs,array,calc}
\usepackage[normalem]{ulem}
\usepackage{xcolor,hyperref}
\hypersetup{hidelinks}
\setcounter{secnumdepth}{-1}
\newcounter{none}
\providecommand{\tightlist}{\setlength{\itemsep}{0pt}\setlength{\parskip}{0pt}}
\providecommand{\passthrough}[1]{#1}
\usepackage{amssymb,mathrsfs,esint,bm,newunicodechar}
\allowdisplaybreaks
\setlength{\emergencystretch}{3em}
\providecommand{\llangle}{\langle\!\langle}
\providecommand{\rrangle}{\rangle\!\rangle}
\newunicodechar{∞}{\ensuremath{\infty}}
\newunicodechar{∈}{\ensuremath{\in}}
\newunicodechar{∑}{\ensuremath{\sum}}
\newunicodechar{∫}{\ensuremath{\int}}
\newunicodechar{≤}{\ensuremath{\leq}}
\newunicodechar{≥}{\ensuremath{\geq}}
\newunicodechar{≠}{\ensuremath{\neq}}
\newunicodechar{→}{\ensuremath{\to}}
\newunicodechar{⇒}{\ensuremath{\Rightarrow}}
\newunicodechar{α}{\ensuremath{\alpha}}
\newunicodechar{β}{\ensuremath{\beta}}
\newunicodechar{γ}{\ensuremath{\gamma}}
\newunicodechar{λ}{\ensuremath{\lambda}}
\newunicodechar{φ}{\ensuremath{\varphi}}
\newunicodechar{ε}{\ensuremath{\varepsilon}}
\newunicodechar{σ}{\ensuremath{\sigma}}
\newunicodechar{θ}{\ensuremath{\theta}}
\newunicodechar{δ}{\ensuremath{\delta}}
\newunicodechar{η}{\ensuremath{\eta}}
\newunicodechar{π}{\ensuremath{\pi}}
\newunicodechar{ξ}{\ensuremath{\xi}}
\newunicodechar{Σ}{\ensuremath{\Sigma}}
\newunicodechar{Δ}{\ensuremath{\Delta}}
\newunicodechar{∂}{\ensuremath{\partial}}
\newunicodechar{∀}{\ensuremath{\forall}}
\newunicodechar{∃}{\ensuremath{\exists}}
\newunicodechar{√}{\ensuremath{\surd}}
\newunicodechar{−}{\ensuremath{-}}
\newunicodechar{μ}{\ensuremath{\mu}}
\newunicodechar{ρ}{\ensuremath{\rho}}
\newunicodechar{τ}{\ensuremath{\tau}}
\newunicodechar{ω}{\ensuremath{\omega}}
\newunicodechar{Ω}{\ensuremath{\Omega}}
\newunicodechar{χ}{\ensuremath{\chi}}
\newunicodechar{ψ}{\ensuremath{\psi}}
\newunicodechar{Γ}{\ensuremath{\Gamma}}
\newunicodechar{Φ}{\ensuremath{\Phi}}
\newunicodechar{Λ}{\ensuremath{\Lambda}}
\newunicodechar{ℝ}{\ensuremath{\mathbb{R}}}
\newunicodechar{ℤ}{\ensuremath{\mathbb{Z}}}
\newunicodechar{ℕ}{\ensuremath{\mathbb{N}}}
\newunicodechar{△}{\ensuremath{\triangle}}
\newunicodechar{⊗}{\ensuremath{\otimes}}
\begin{document}
\section{2025年杭州电子科技大学601数学分析考研真题}\label{ux5e74ux676dux5ddeux7535ux5b50ux79d1ux6280ux5927ux5b66601ux6570ux5b66ux5206ux6790ux8003ux7814ux771fux9898}

来源：\texttt{HDU/601数学分析/2025/2025年杭州电子科技大学601数学分析考研真题.pdf}。原
PDF 页序：1, 2。

\section{杭州电子科技大学2025年数学分析真题}\label{ux676dux5ddeux7535ux5b50ux79d1ux6280ux5927ux5b662025ux5e74ux6570ux5b66ux5206ux6790ux771fux9898}

\begin{quote}
文字整理说明：根据远程OCR与本地逐页OCR整理数学语法，一6的积分式由本地OCR补全。保留原资料的缺失条件和待讨论结论，不将补充解析中的建议条件写入题干。
\end{quote}

\subsection{一、填空（每题5分，共40分）}\label{ux4e00ux586bux7a7aux6bcfux98985ux5206ux517140ux5206}

\begin{enumerate}
\def\labelenumi{\arabic{enumi}.}
\item
  \(\displaystyle\lim_{n\to+\infty}\frac{11\times12\times\cdots\times(n+10)}{1\times4\times7\times\cdots\times(3n-2)}=\)
  \_\_\_\_。
\item
  \(\forall x>0\)，\(f(x)>0\)；\(\forall x\ge0\)，\(\displaystyle f(x)=\sqrt{\int_0^xf(t)\,dt}\)，则
  \(f(x)=\) \_\_\_\_。
\item
  \(\displaystyle\int_\alpha^{\alpha+1}(x-1)(1-x)^{10}\,dx=\) \_\_\_\_。
\item
  若
  \(\displaystyle\lim_{x\to0}(1+x)^{c/x}=\int_{-\infty}^cte^t\,dt\)，则
  \(c=\) \_\_\_\_。
\item
  \(\displaystyle\iiint_V(x+y+z)\,dx\,dy\,dz=\) \_\_\_\_。（\(V\) 为
  \(x+y+z=1\)
  和三个坐标轴围成。原资料用词如此，几何含义的说明见补充解析。）
\item
  \(\displaystyle\int_Lx\,dy=\) \_\_\_\_。（\(L\) 是
  \(x^2/4+y^2=1\)，逆时针。）
\item
  隐函数中拐点，求参数。（原资料未给出具体方程和参数条件。）
\item
  \(\displaystyle\iint_S(x^2+y^2)\,dS=\) \_\_\_\_，\(S\) 为
  \(z=\sqrt{x^2+y^2}\) 和 \(z=1\) 所围的外表面。
\end{enumerate}

\subsection{二、计算（每题9分，共36分）}\label{ux4e8cux8ba1ux7b97ux6bcfux98989ux5206ux517136ux5206}

\begin{enumerate}
\def\labelenumi{\arabic{enumi}.}
\item
  计算 \(\displaystyle\iint_Sx^2y^2z\,dx\,dy\)，\(S\) 为
  \(x^2+y^2+z^2=1\) 的下半球面，取上侧。
\item
  \(\displaystyle I=\int_LP\,dx+Q\,dy\)，其中 \[
  P=\frac{x}{y}r^\alpha,\quad Q=-\frac{x^2}{y^2}r^\alpha,\quad r=\sqrt{x^2+y^2}.
  \] 若 \(I\equiv0\)，则 \(\alpha\) 为多少？（现有题面未给出 \(L\)
  的定义。）
\item
  求 \(\displaystyle\sum_{n=2}^\infty\frac{(-1)^nx^{n+1}}{n(n-1)}\)
  的收敛半径和和函数。
\item
  若 \(\displaystyle\int_0^\infty e^{-x^2}\,dx=\frac{\sqrt\pi}{2}\)，求
  \(\displaystyle\int_0^\infty\cos(\alpha x)e^{-x^2}\,dx\)。
\end{enumerate}

\subsection{三、证明题（每题10分，共50分）}\label{ux4e09ux8bc1ux660eux9898ux6bcfux989810ux5206ux517150ux5206}

\begin{enumerate}
\def\labelenumi{\arabic{enumi}.}
\item
  用``\(\varepsilon\)---\(\delta\)''语言证明： \[
  \lim_{x\to1}\frac{x^2-x}{6x^2-4x-2}=\frac18.
  \]
\item
  \(f(x)\)
  二阶连续可导，\(f(0)=0\)，\(\displaystyle M=\max_{x\in[0,1]}f''(x)\)，证明：
  \[
  \left|\int_{-1}^1f(x)\,dx\right|\le\frac13M.
  \] （\(M\)
  的区间及未加绝对值的写法均按现有题面保留；命题的条件问题见补充解析。）
\item
  证明： \[
  u(x,t)=\frac12[f(x-at)+f(x+at)]+\frac1{2a}\int_{x-at}^{x+at}F(z)\,dz
  \] 满足弦振动方程
  \(\displaystyle\frac{\partial^2u}{\partial t^2}=a^2\frac{\partial^2u}{\partial x^2}\)。
\item
  证明函数列
  \(\displaystyle S_n(x)=\sum_{k=1}^n\frac1n\cos\left(x+\frac kn\right)\)
  在 \(x\in[a,b]\subset\mathbb R\) 上一致收敛。
\item
  若 \(f_n(x)\) 在区间 \(I\) 上一致连续，且
  \(f_n\rightrightarrows f\)，证明 \(f\) 在 \(I\) 上一致连续。
\end{enumerate}

\subsection{四、讨论题（每题12分，共24分）}\label{ux56dbux8ba8ux8bbaux9898ux6bcfux989812ux5206ux517124ux5206}

\begin{enumerate}
\def\labelenumi{\arabic{enumi}.}
\item
  若
  \(\displaystyle\frac{x^2}{\alpha^2}+\frac{y^2}{\beta^2}+\frac{z^2}{\gamma^2}=1\)，切平面与坐标轴围成的四面体体积为
  \(V\)，求
  \(V_{\max}\)。（``最大值''按原资料保留，存在性问题见补充解析。）
\item
  讨论函数 \[
  f(x)=\begin{cases}x(1-x),&x\in\mathbb Q,\\x(1+x),&x\notin\mathbb Q\end{cases}
  \] 的连续性和可微性。
\end{enumerate}
\end{document}
